
\vspace{1cm}

\noindent  \textbf{Leia as instruções antes de fazer a lista:}

\vspace{0.3cm}
\noindent A resolução deve ser entregue em formato Jupyter Notebook ou Google Colab na plataforma Google Classroom, na tarefa designada e até o dia determinado.

\noindent  Não é necessário testar se os dados passados por argumento são válidos a menos que pedido na questão.


\vspace{0.5cm}


\begin{enumerate}[label=\textbf{Q\arabic*.}]

% ---------------------------------------------------------
\item Em uma estação meteorológica, os dados diários de temperatura são gravados em uma tupla de valores numéricos. Crie uma função chamada \code{resumir\_temperaturas(temperaturas)} que extraia, via fatiamento (\emph{slicing}), apenas as leituras intermediárias (descartando a primeira e a última medição) e retorne três valores usando retorno múltiplo: a menor temperatura do período intermediário, a maior e a média aritmética das temperaturas intermediárias.

\begin{itemize}[leftmargin=*]
    \item \textbf{Entrada:} \code{(18.0, 21.5, 25.0, 23.0, 19.5, 17.0)}
    \item \textbf{Saída:} \code{(19.5, 25.0, 22.25)}
\end{itemize}


\ifmostrarGabarito
\begin{minted}{python}
# Resposta Questao 1
def resumir_temperaturas(temperaturas):
    # Fatiamento descarta o primeiro [0] e o ultimo [-1] elementos
    intermediarias = temperaturas[1:-1]
    
    menor = min(intermediarias)
    maior = max(intermediarias)
    media = sum(intermediarias) / len(intermediarias)
    
    return menor, maior, media
\end{minted}
\fi




\item Uma biblioteca comunitária está reorganizando seu catálogo digital. Os títulos dos livros foram inseridos com traços baixos (\textit{underline}) substituindo os espaços e uma etiqueta no final: \code{"O\_Alquimista-ROMANCE"}. Escreva uma função \code{formatar\_catalogo(lista\_raw)} que receba uma lista de strings nesse formato, utilize \code{.split('-')} para desempacotar o título bruto e a categoria, substitua os traços baixos (\code{\_}) por espaços no título usando \code{.replace()}, e retorne uma lista de tuplas \code{[(titulo\_limpo, categoria)]}.

\begin{itemize}[leftmargin=*]
    \item \textbf{Entrada:} \code{["O\_Alquimista-ROMANCE", "Dom\_Casmurro-CLASSICO"]}
    \item \textbf{Saída:} \code{[('O Alquimista', 'ROMANCE'), ('Dom Casmurro', 'CLASSICO')]}
\end{itemize}

\ifmostrarGabarito{
\begin{minted}{python}
# Resposta Questao 2
def formatar_catalogo(lista_raw):
    catalogo = []
    for item in lista_raw:
        titulo_bruto, categoria = item.split("-")
        titulo_limpo = titulo_bruto.replace("_", " ")
        catalogo.append((titulo_limpo, categoria))
    return catalogo
\end{minted}
}



\item Um aplicativo de corridas urbanas registra o histórico de viagens de um motorista em uma lista de tuplas contendo \code{(id\_corrida, valor\_reais, (latitude\_destino, longitude\_destino))}. Crie uma função chamada \code{filtrar\_corridas\_rentaveis(historico, valor\_minimo)} que percorra a lista e retorne uma lista de dicionários contendo apenas as corridas cujo valor seja estritamente superior ao \code{valor\_minimo}. O dicionário deve possuir as chaves \code{"id"} e \code{"destino"}.

\begin{itemize}[leftmargin=*]
    \item \textbf{Entrada:} \code{[(101, 25.0, (-23.55, -46.63)), (102, 12.0, (-23.58, -46.65))], 15.0}
    \item \textbf{Saída:} \code{[\{'id': 101, 'destino': (-23.55, -46.63)\}]}
\end{itemize}

\ifmostrarGabarito
\begin{minted}{python}
# Resposta Questao 3
def filtrar_corridas_rentaveis(historico, valor_minimo):
    rentaveis = []
    for id_corrida, valor, destino in historico:
        if valor > valor_minimo:
            rentaveis.append({
                "id": id_corrida,
                "destino": destino
            })
    return rentaveis
\end{minted}
\fi



\item Em um sistema de cadastro de clientes de um banco, passar uma lista mutável para uma função pode alterar a lista original indesejadamente. 

Crie uma função \code{adicionar\_tag\_segura(lista\_tags, nova\_tag)} que receba uma lista de tags de um cliente. A função deve criar uma cópia rasa (\emph{shallow copy}), utilizando o método \code{.copy()} da lista antes de adicionar a \code{nova\_tag}, garantindo que a lista original passada como argumento não seja modificada. O retorno deve ser a nova lista.

\begin{itemize}[leftmargin=*]
    \item \textbf{Entrada:} \code{["VIP", "Investidor"]}, \code{"Premium"}
    \item \textbf{Saída (retorno da função):} \code{['VIP', 'Investidor', 'Premium']}
    \item \textbf{Verificação da lista original:} Permanece \code{['VIP', 'Investidor']}
\end{itemize}
\ifmostrarGabarito{
\begin{minted}{python}
# Resposta Questao 4
def adicionar_tag_segura(lista_tags, nova_tag):
    # Cria uma copia da lista para evitar mutacao de lado na original
    copia_lista = lista_tags.copy()
    copia_lista.append(nova_tag)
    return copia_lista
\end{minted}
}









\item Em uma plataforma de e-commerce de artigos esportivos, os dados de venda chegam em uma string consolidada onde as compras são separadas por ponto e vírgula (\code{;}) e os atributos do item por cerquilha (\code{\#}): \code{"Bola\#3\#45.0;Chuteira\#1\#180.0"}. Escreva uma função chamada \code{consolidar\_carrinho(texto\_compras)} que processe a string e retorne uma lista de tuplas no formato \code{[(produto\_str, quantidade\_int, subtotal\_float)]}.

\begin{itemize}[leftmargin=*]
    \item \textbf{Entrada:} \code{"Bola\#3\#45.0;Chuteira\#1\#180.0"}
    \item \textbf{Saída:} \code{[('Bola', 3, 135.0), ('Chuteira', 1, 180.0)]}
\end{itemize}

\ifmostrarGabarito
\begin{minted}{python}
# Resposta Questao 5
def consolidar_carrinho(texto_compras):
    carrinho = []
    itens = texto_compras.split(";")
    for item in itens:
        produto, qtd_str, preco_str = item.split("#")
        qtd = int(qtd_str)
        subtotal = qtd * float(preco_str)
        carrinho.append((produto, qtd, subtotal))
    return carrinho
\end{minted}
\fi



Uma escola de música organiza suas turmas em uma estrutura integrada. Ela possui um cadastro de alunos no formato de lista de dicionários:

\begin{minted}{python}
alunos = [
    {"nome": "Lucas", "instrumento": "Violao", "notas": "8.5,9.0,7.5"},
    {"nome": "Beatriz", "instrumento": "Piano", "notas": "10.0,9.5,9.0"}
]
\end{minted}

Escreva uma função chamada \code{gerar\_relatorio\_musica(lista\_alunos)} que processe esses dados e retorne um novo dicionário onde:
\begin{itemize}[leftmargin=*]
    \item A \textbf{chave} seja uma tupla imutável com o nome do aluno e seu instrumento: \code{(nome, instrumento)}.
    \item O \textbf{valor} seja uma tupla contendo a média das notas (arredondada para 2 casas decimais) e uma string de status (\code{"Aprovado"} se média $\ge 8.0$, senão \code{"Em Desenvolvimento"}).
\end{itemize}

\begin{itemize}[leftmargin=*]
    \item \textbf{Entrada:} A lista \code{alunos} fornecida acima.
    \item \textbf{Saída:} \code{\{('Lucas', 'Violao'): (8.33, 'Aprovado'),}\\
    \code{('Beatriz', 'Piano'): (9.5, 'Aprovado')\}}
\end{itemize}

\ifmostrarGabarito
\begin{minted}{python}
# Resposta Questão 6
def gerar_relatorio_colecao(garagem):
    relatorio = {}
    
    for carro in garagem:
        # Cria a tupla imutável (modelo, ano) para 
        # ser usada como chave do dicionário
        chave_carro = (carro["modelo"], carro["ano"])
        
        # Separa a string de donos pela vírgula e 
        # remove os espaços extras de cada nome
        lista_donos_raw = carro["donos"].split(",")
        lista_donos_limpa = [dono.strip() for dono in lista_donos_raw]
        
        # Atribui a lista de donos processada à chave correspondente
        relatorio[chave_carro] = lista_donos_limpa
        
    return relatorio
\end{minted}
\fi



\end{enumerate}